\section{Introduction}
\label{sec:introduction}

Deep learning for financial time series has accumulated a strong empirical
record on single-market benchmarks, but the question of whether a
\emph{learned architecture} transfers across structurally different
markets remains largely open. In limit order book (LOB) microstructure,
convolutional and recurrent backbones such as DeepLOB
\cite{zhang2019deeplob} and its descendants \cite{briola2021deep,
lit2025limit} routinely outperform tabular baselines.  In decentralized
finance (DeFi) lending, by contrast, the published literature has so far
favored reinforcement learning either on the protocol side
\cite{bertucci2025rl, xu2023autogov} or recursive least-squares adaptive
control \cite{agilerate2024}.  A practitioner who has invested in a
domain-aware recurrent stack for one of these markets faces an honest
question: must the model be re-engineered from scratch for the other, or
is the architectural template itself portable?

The two markets are not obviously compatible. LOB mid-price prediction
operates on \mbox{millisecond-scale} order arrivals with a discrete
tick lattice and a strong micro\-structural informational asymmetry.
DeFi lending rates, by contrast, are set deterministically from on-chain
state via a known piecewise-linear kink function
$f_{\text{kink}}(u)$ of utilization $u$; they update hourly at most and
move primarily because of cointegration between protocols
\cite{gudgeon2020loanable} and slow demand-driven re-pricing
\cite{baude2025optimal}.  The two are noisy multivariate time series
with regime-dependent dynamics and a weighted-error preference for large
moves, but their feature topologies, decision horizons, and competitor
adversaries differ qualitatively.

\paragraph{Central thesis.}
We test the conjecture that \emph{the same dual-branch architecture
trained with the same composite loss generalizes between LOB mid-price
prediction and DeFi lending-rate forecasting by re-identifying the
domain-natural decomposition pair.}  In LOB the pair is (price, volume)
augmented with 21 microstructure features
\cite{solovev2026lob}; in DeFi it is
(rate residual $\varepsilon = r - f_{\text{kink}}(u)$,
utilization with context).  The two BiGRU branches, multi-scale Conv1d
fusion head, and the composite loss
$\alpha\,\textsc{MSE} + \beta\,(1{-}\rho_w) + \gamma\,Q_{90}$ are
unchanged between the two specialisations; only the input subspaces and
their dimensionality are re-mapped to the target domain.

\paragraph{Contributions.}
The paper makes four contributions.
\begin{itemize}
  \item \textbf{A unified dual-branch recurrent architecture
    (DA-BiGRU-CNN) tested in two structurally different financial domains
    --- LOB mid-price and DeFi USDC lending rates --- under the same
    composite loss.}
  \item \textbf{A new empirical study: to our knowledge the first
    published forecast-driven multi-criteria allocator across Aave V3
    and Compound V3 USDC pools}, trained on $\DataHours{}$ hourly bars
    spanning \DataStart{}--\DataEnd{} at $\DataCoverage$ panel coverage,
    with regime structure documented at full coverage.
  \item \textbf{An honest negative-result protocol.}  We pre-register the
    headline H1 ($\Delta\textsc{Sharpe} \ge 0.2$) and report
    $\Delta\textsc{Sharpe} = \SharpeDelta$ with the bootstrap
    95\% CI $[\BootstrapCILow,\,\BootstrapCIHigh]$ and $p$-value
    $\BootstrapPvalue$ --- failing to reject H0 under the binding
    $n_{\text{months}}=4$ test window, exactly as the pre-registration
    permitted.
  \item \textbf{Two upstream pull requests to the open-source
    \texttt{fractal-defi} framework} \cite{krestenko2026fractaldefi}: a
    Compound V3 Messari loader and a \texttt{Base\-Lending\-Allocation\-Strategy}
    abstraction with cooldown infrastructure --- both released alongside
    the paper for reproducibility.
\end{itemize}

\paragraph{Outline.}
\Cref{sec:background} reviews LOB microstructure DL, DeFi lending
mechanics, the impossibility result of \cite{halpern2024fair}, and
related transfer-learning literature.  \Cref{sec:methodology} gives the
domain-agnostic architecture and its two specialisations.
\Cref{sec:lob-recap} summarises the LOB results we treat as prior work.
\Cref{sec:defi-experiment} reports the DeFi empirical study with data,
forecaster training, MCDM allocator, headline H1 test, and ablations.
\Cref{sec:discussion} states what transfers and what does not, and
\Cref{sec:limitations,sec:conclusion} discuss honest limitations and
close.
